\documentclass[12pt,a4paper,oneside]{report}

\usepackage[french]{babel}
\usepackage[T1]{fontenc}
\usepackage{geometry}
\usepackage{fancyhdr}
\usepackage{graphicx}
\usepackage{listings}
\usepackage{xcolor}
\usepackage{amsmath}
\usepackage{amssymb}
\usepackage{booktabs}
\usepackage{array}
\usepackage{multirow}
\usepackage{caption}
\usepackage{float}
\usepackage{tocloft}
\usepackage{tikz}
\usepackage{pgfplots}
\usepackage{algorithm}
\usepackage{algpseudocode}
\usepackage{hyperref}
\usepackage{setspace}
\usepackage[utf-8]{inputenc}

% Geometry
\geometry{
  left=2.5cm,
  right=2.5cm,
  top=2.5cm,
  bottom=2.5cm
}

% Code listing configuration
\lstset{
  basicstyle=\ttfamily\small,
  breaklines=true,
  breakatwhitespace=true,
  commentstyle=\color{gray},
  keywordstyle=\color{blue},
  numberstyle=\tiny\color{gray},
  stringstyle=\color{darkgreen},
  showstringspaces=false,
  tabsize=2,
  language=C++,
  morekeywords={rclcpp, Node, create_subscription, create_publisher, std_msgs, geometry_msgs},
  backgroundcolor=\color{lightgray!10},
  frame=single,
  lineskip=-1pt
}

% Header and Footer
\pagestyle{fancy}
\fancyhf{}
\rhead{\thepage}
\lhead{Système ROS2 - Projet Robot Autonome}
\renewcommand{\headrulewidth}{0.4pt}

% Color definitions
\definecolor{darkgreen}{rgb}{0,0.5,0}
\definecolor{darkred}{rgb}{0.5,0,0}
\definecolor{lightgray}{rgb}{0.9,0.9,0.9}

% Title Page
\begin{document}

\begin{titlepage}
  \centering
  \vspace*{2cm}
  {\LARGE\textbf{Rapport du Projet de Fin d'Études}}
  \\[0.5cm]
  {\Large\textbf{Architecture et Implémentation d'un Système ROS2 pour Exploration Robotique Autonome}}
  \\[1cm]
  \vspace{2cm}
  
  \begin{flushleft}
    \textbf{Titre du projet:} Système de Navigation et d'Exploration Autonome pour Robot Mobile \\[0.5cm]
    \textbf{Type de robot:} TurtleBot-like Mobile Robot \\[0.5cm]
    \textbf{Cadre de travail:} Robot Operating System 2 (ROS2) \\[0.5cm]
    \textbf{Date de soumission:} \today
  \end{flushleft}
  
  \vfill
  
  \begin{center}
    {\large\textbf{Résumé Exécutif}}
    \\[0.5cm]
    Ce rapport présente une analyse détaillée et complète de l'architecture d'un système ROS2 conçu pour l'exploration autonome d'environnements inconnus par un robot mobile. Le projet intègre des composants sophistiqués de détection des frontières, de coordination de l'exploration, de navigation, et de sécurité. L'architecture repose sur une communication par sujets ROS2 (topics) et des actions (actions), permettant une coordination efficace entre les différents nœuds du système.
  \end{center}
  
  \vspace{1cm}
  
\end{titlepage}

\newpage
\tableofcontents
\newpage

\chapter{Introduction et Contexte}

\section{Objectifs du Projet}

Le projet décrit dans ce rapport vise à développer un système complètement autonome pour l'exploration d'environnements inconnus. Les objectifs principaux sont:

\begin{enumerate}
  \item Implémenter un système de détection des frontières permettant d'identifier les zones inexploitées
  \item Créer un coordinateur d'exploration capable de gérer plusieurs états et transitions
  \item Assurer la sécurité du robot à travers des mécanismes d'arrêt d'urgence et de filtrage de vitesse
  \item Intégrer la localisation et la cartographie simultanées (SLAM)
  \item Fournir des outils de visualisation en temps réel pour le suivi de l'exploration
\end{enumerate}

\section{Architecture Générale du Système}

Le système ROS2 est composé de plusieurs paquets interdépendants:

\begin{itemize}
  \item \textbf{Bringup:} Responsable du lancement et de l'initialisation de tous les composants
  \item \textbf{Exploration:} Contient la détection des frontières et la coordination de l'exploration
  \item \textbf{Navigation:} Gère la navigation, l'arrêt d'urgence et le suivi des murs
  \item \textbf{SLAM:} Effectue la localisation et la cartographie simultanées
  \item \textbf{Description:} Fournit la description URDF du robot et les contrôleurs
  \item \textbf{Gazebo:} Simule l'environnement et la physique du robot
\end{itemize}

\chapter{Fondamentaux ROS2 et Architecture Distribuée}

\section{Introduction à ROS2}

ROS2 (Robot Operating System 2) est un framework distribué conçu pour faciliter le développement de systèmes robotiques complexes. Contrairement à ROS1, ROS2 utilise le middleware DDS (Data Distribution Service) comme couche de communication, offrant une meilleure scalabilité, une latence réduite et une meilleure intégration avec les systèmes temps réel.

\subsection{Concepts Fondamentaux}

\subsubsection{Nœuds (Nodes)}

Un nœud est une unité exécutable autonome dans ROS2. Chaque nœud réalise une fonction spécifique et communique avec d'autres nœuds. Dans notre système, nous avons plusieurs nœuds critiques:

\begin{itemize}
  \item \textbf{Frontier Detection Node:} Detecte les frontières dans la carte
  \item \textbf{Exploration Coordinator Node:} Coordonne la stratégie d'exploration
  \item \textbf{Emergency Stop Controller Node:} Assure la sécurité
  \item \textbf{Wall Follower Node:} Implémente le suivi de mur comme stratégie alternative
  \item \textbf{Navigation Stack Nodes:} Composants Navigation2
  \item \textbf{SLAM Nodes:} Effectue la cartographie
\end{itemize}

\subsubsection{Sujets (Topics)}

Les sujets (topics) sont des canaux de communication asynchrones permettant la publication et l'abonnement de messages. La communication par sujets utilise un modèle publish-subscribe où les producteurs de données publient sur un sujet et les consommateurs s'abonnent au sujet.

\subsubsection{Services (Services)}

Les services fournissent une communication synchrone de requête-réponse entre les nœuds.

\subsubsection{Actions (Actions)}

Les actions permettent une communication asynchrone avec feedback, idéale pour les tâches longues comme la navigation vers un objectif.

\section{Modèle de Communication}

Notre système utilise un modèle de communication distribué basé sur les topics et les actions. La figure suivante illustre le flux de données:

\begin{center}
\begin{tikzpicture}[node distance=2cm and 3cm, auto]
  \node [draw, rectangle] (slam) {SLAM Node};
  \node [draw, rectangle, right of=slam] (fd) {Frontier Detector};
  \node [draw, rectangle, below of=fd] (ec) {Exploration Coordinator};
  \node [draw, rectangle, left of=ec] (nav2) {Navigation2};
  \node [draw, rectangle, right of=ec] (esc) {Emergency Stop};
  
  \draw [->, thick] (slam) -- (fd) node [midway, above] {/map};
  \draw [->, thick] (fd) -- (ec) node [midway, right] {/exploration/frontiers};
  \draw [->, thick] (ec) -- (nav2) node [midway, below] {navigate\_to\_pose};
  \draw [->, thick] (ec) -- (esc) node [midway, right] {/cmd\_vel};
  \draw [->, thick] (esc) -- (nav2) node [midway, below] {/cmd\_vel\_safe};
  
\end{tikzpicture}
\end{center}

\chapter{Système de Détection des Frontières}

\section{Vue d'ensemble}

Le système de détection des frontières est responsable d'identifier les zones de transition entre l'espace exploré et l'espace inexploré dans la carte. Une frontière est définie comme un ensemble de cellules libres adjacentes à des cellules inconnues.

\section{Implémentation du Nœud Frontier Detector}

\subsection{Initialisation et Configuration}

Le nœud Frontier Detector hérite de la classe \texttt{rclcpp::Node} et initialise plusieurs paramètres configurables:

\begin{lstlisting}
FrontierDetector::FrontierDetector(const rclcpp::NodeOptions & options)
: Node("frontier_detector", options), map_received_(false)
{
  // Declare parameters with default values
  this->declare_parameter("min_frontier_size", 10.0);
  this->declare_parameter("max_frontier_distance", 10.0);
  this->declare_parameter("frontier_travel_point_distance", 0.5);
  this->declare_parameter("information_radius", 1.0);
  this->declare_parameter("robot_base_frame", "base_link");
  this->declare_parameter("global_frame", "map");
  
  // Get parameters from ROS2 parameter server
  min_frontier_size_ = 
    this->get_parameter("min_frontier_size").as_double();
  max_frontier_distance_ = 
    this->get_parameter("max_frontier_distance").as_double();
  
  // Initialize TF2 for coordinate transformations
  tf_buffer_ = std::make_shared<tf2_ros::Buffer>(
    this->get_clock());
  tf_listener_ = 
    std::make_shared<tf2_ros::TransformListener>(*tf_buffer_);
}
\end{lstlisting}

\subsection{Abonnements (Subscriptions) et Publications (Publishers)}

Le nœud Frontier Detector s'abonne à la carte occupancy grid et publie les frontières détectées ainsi que des marqueurs de visualisation:

\begin{lstlisting}
// Subscription to occupancy grid map
map_sub_ = this->create_subscription<nav_msgs::msg::OccupancyGrid>(
  "/map", rclcpp::QoS(1).transient_local(),
  std::bind(&FrontierDetector::mapCallback, this, 
    std::placeholders::_1));

// Publisher for frontier visualization
frontier_marker_pub_ = 
  this->create_publisher<visualization_msgs::msg::MarkerArray>(
    "/exploration/frontiers", 10);

// Publisher for frontier goals
frontier_goal_pub_ = 
  this->create_publisher<geometry_msgs::msg::PoseStamped>(
    "/exploration/frontier_goal", 10);

// Timer for periodic detection (1000ms interval)
timer_ = this->create_wall_timer(
  std::chrono::milliseconds(1000),
  std::bind(&FrontierDetector::timerCallback, this));
\end{lstlisting}

\subsection{Algorithme de Détection des Frontières}

L'algorithme de détection fonctionne en plusieurs étapes:

\begin{enumerate}
  \item \textbf{Conversion d'image:} Convertir la grille d'occupation en image OpenCV
  \item \textbf{Identification des cellules frontière:} Trouver toutes les cellules libres adjacentes à des cellules inconnues
  \item \textbf{Groupement:} Regrouper les cellules frontière en frontières distinctes
  \item \textbf{Filtrage:} Filtrer par taille minimale et distance maximale au robot
  \item \textbf{Calcul du gain informatif:} Évaluer la valeur d'exploration de chaque frontière
  \item \textbf{Vérification d'accessibilité:} Assurer que la frontière est atteignable
\end{enumerate}

\begin{lstlisting}
std::vector<Frontier> FrontierDetector::detectFrontiers(
  const nav_msgs::msg::OccupancyGrid & map)
{
  // Convert occupancy grid to OpenCV image format
  cv::Mat map_image(map.info.height, map.info.width, CV_8UC1);
  
  for (unsigned int y = 0; y < map.info.height; ++y) {
    for (unsigned int x = 0; x < map.info.width; ++x) {
      int index = y * map.info.width + x;
      int8_t value = map.data[index];
      
      if (value == -1) {
        map_image.at<uint8_t>(y, x) = 128;  // Unknown
      } else if (value == 0) {
        map_image.at<uint8_t>(y, x) = 255;  // Free
      } else {
        map_image.at<uint8_t>(y, x) = 0;    // Occupied
      }
    }
  }
  
  // Find frontier cells and group them
  auto frontier_cells = findFrontierCells(map_image);
  auto frontiers = groupFrontierCells(frontier_cells, map);
  
  // Filter frontiers by criteria
  std::vector<Frontier> filtered_frontiers;
  auto robot_pose = getRobotPose();
  
  for (auto & frontier : frontiers) {
    if (frontier.size >= min_frontier_size_) {
      // Calculate distance to robot
      double dx = frontier.centroid.x - robot_pose.pose.position.x;
      double dy = frontier.centroid.y - robot_pose.pose.position.y;
      frontier.distance_to_robot = std::sqrt(dx * dx + dy * dy);
      
      if (frontier.distance_to_robot <= max_frontier_distance_) {
        frontier.information_gain = 
          calculateInformationGain(frontier, map);
        frontier.is_reachable = isReachable(frontier.centroid);
        
        if (frontier.is_reachable) {
          filtered_frontiers.push_back(frontier);
        }
      }
    }
  }
  
  return filtered_frontiers;
}
\end{lstlisting}

\subsection{Identification des Cellules Frontière}

Une cellule frontière est une cellule libre ayant au moins un voisin inconnu:

\begin{lstlisting}
bool FrontierDetector::isFrontierCell(
  const cv::Mat & map, int x, int y)
{
  // A frontier cell is a free cell with at least one unknown neighbor
  if (map.at<uint8_t>(y, x) != 255) {  // Not free space
    return false;
  }
  
  // Check 8-neighbor connectivity
  for (int dy = -1; dy <= 1; ++dy) {
    for (int dx = -1; dx <= 1; ++dx) {
      if (dx == 0 && dy == 0) continue;
      
      int nx = x + dx;
      int ny = y + dy;
      
      if (nx >= 0 && nx < map.cols && ny >= 0 && ny < map.rows) {
        if (map.at<uint8_t>(ny, nx) == 128) {  // Unknown neighbor
          return true;
        }
      }
    }
  }
  
  return false;
}
\end{lstlisting}

\section{Publications et Sujets}

Les frontières détectées sont publiées sous forme de marqueurs RViz pour la visualisation:

\begin{lstlisting}
// Topic: /exploration/frontiers
// Type: visualization_msgs/MarkerArray
// Contient les positions et les gains informatifs de chaque frontière

// Topic: /exploration/frontier_goal
// Type: geometry_msgs/PoseStamped
// Publie l'objectif de frontière sélectionné
\end{lstlisting}

\chapter{Coordinateur d'Exploration}

\section{Vue d'ensemble}

Le Exploration Coordinator est le cerveau du système d'exploration. Il maintient un état machine (state machine) et orchestre l'exploration en sélectionnant les frontières optimales et en commandant le robot pour les atteindre.

\section{Machine à États}

Le coordinateur gère plusieurs états distincts:

\begin{enumerate}
  \item \textbf{IDLE:} Système au repos, attendant l'activation
  \item \textbf{DETECTING\_ENVIRONMENT:} Phase de détection initiale de l'environnement
  \item \textbf{EXPLORING:} Exploration active vers une frontière
  \item \textbf{SELECTING\_GOAL:} Sélection de la prochaine frontière cible
  \item \textbf{MOVING\_TO\_GOAL:} Navigation vers l'objectif
  \item \textbf{GOAL\_REACHED:} Objectif atteint
  \item \textbf{STUCK:} Robot bloqué, nécessite récupération
  \item \textbf{RECOVERY:} Mode récupération actif
\end{enumerate}

\section{Initialisation et Paramètres}

\begin{lstlisting}
ExplorationCoordinator::ExplorationCoordinator(
  const rclcpp::NodeOptions & options)
: Node("exploration_coordinator", options),
  current_state_(ExplorationState::IDLE),
  map_received_(false),
  frontiers_received_(false),
  exploration_active_(false)
{
  // Declare exploration parameters
  this->declare_parameter("exploration_timeout", 300.0);
  this->declare_parameter("stuck_timeout", 30.0);
  this->declare_parameter("min_map_coverage_to_start", 0.1);
  this->declare_parameter("goal_tolerance", 0.8);
  this->declare_parameter("stuck_distance_threshold", 0.15);
  this->declare_parameter("goal_reached_tolerance", 1.2);
  this->declare_parameter("goal_timeout", 45.0);
  this->declare_parameter("force_new_goal_distance", 1.5);
  
  // Declare optimization parameters
  this->declare_parameter("enable_dynamic_tolerance", true);
  this->declare_parameter("velocity_based_tolerance", true);
  this->declare_parameter("environmental_adaptation", true);
  this->declare_parameter("performance_monitoring", true);
  this->declare_parameter("adaptive_parameter_tuning", true);
  
  // Initialize TF2
  tf_buffer_ = std::make_shared<tf2_ros::Buffer>(
    this->get_clock());
  tf_listener_ = std::make_shared<tf2_ros::TransformListener>(
    *tf_buffer_);
}
\end{lstlisting}

\section{Abonnements aux Topics}

Le coordinateur s'abonne aux topics critiques:

\begin{lstlisting}
// Subscribe to occupancy grid
map_sub_ = this->create_subscription<nav_msgs::msg::OccupancyGrid>(
  "/map", rclcpp::QoS(1).transient_local(),
  std::bind(&ExplorationCoordinator::mapCallback, this, 
    std::placeholders::_1));

// Subscribe to frontier markers
frontier_markers_sub_ = 
  this->create_subscription<visualization_msgs::msg::MarkerArray>(
    "/exploration/frontiers", 10,
    std::bind(&ExplorationCoordinator::frontierMarkersCallback, 
      this, std::placeholders::_1));
\end{lstlisting}

\section{Publications}

\begin{lstlisting}
// Publish exploration status
status_pub_ = this->create_publisher<std_msgs::msg::String>(
  "/exploration/status", 10);

// Publish visualization markers
exploration_markers_pub_ = 
  this->create_publisher<visualization_msgs::msg::MarkerArray>(
    "/exploration/markers", 10);

// Publish velocity commands
cmd_vel_pub_ = this->create_publisher<geometry_msgs::msg::Twist>(
  "/cmd_vel", 10);
\end{lstlisting}

\section{Client d'Action Navigation}

Le coordinateur utilise un client d'action pour naviguer vers les frontières:

\begin{lstlisting}
// Create action client for Navigation2
nav_action_client_ = 
  rclcpp_action::create_client<nav2_msgs::action::NavigateToPose>(
    this, "/navigate_to_pose");

// Send goal to navigation
auto goal = nav2_msgs::action::NavigateToPose::Goal();
goal.pose.header.frame_id = "map";
goal.pose.pose = frontier_pose;

auto send_goal_options = 
  rclcpp_action::Client<nav2_msgs::action::NavigateToPose>::
    SendGoalOptions();
send_goal_options.goal_response_callback = 
  std::bind(&ExplorationCoordinator::goalResponseCallback, this, 
    std::placeholders::_1);
send_goal_options.feedback_callback = 
  std::bind(&ExplorationCoordinator::feedbackCallback, this, 
    std::placeholders::_1, std::placeholders::_2);
send_goal_options.result_callback = 
  std::bind(&ExplorationCoordinator::resultCallback, this, 
    std::placeholders::_1);

nav_action_client_->async_send_goal(goal, send_goal_options);
\end{lstlisting}

\section{Sélection Intelligente de Frontière}

Le coordinateur utilise une stratégie multi-critères pour sélectionner la meilleure frontière:

\begin{lstlisting}
Frontier ExplorationCoordinator::selectBestFrontier(
  const std::vector<Frontier> & frontiers)
{
  double max_score = -std::numeric_limits<double>::infinity();
  Frontier best_frontier;
  
  for (const auto & frontier : frontiers) {
    // Multi-criteria scoring function
    double distance_score = 
      distance_efficiency_weight_ / (1.0 + frontier.distance_to_robot);
    double information_score = 
      frontier.information_gain * exploration_value_weight_;
    double time_score = 
      time_efficiency_weight_ / (1.0 + frontier.distance_to_robot);
    
    double total_score = distance_score + information_score + time_score;
    
    if (total_score > max_score) {
      max_score = total_score;
      best_frontier = frontier;
    }
  }
  
  return best_frontier;
}
\end{lstlisting}

\section{Gestion d'État}

La transition d'état est gérée par une machine à états complète:

\begin{lstlisting}
void ExplorationCoordinator::setState(ExplorationState new_state)
{
  if (new_state == current_state_) {
    return;  // No state change
  }
  
  current_state_ = new_state;
  state_start_time_ = std::chrono::steady_clock::now();
  
  switch (current_state_) {
    case ExplorationState::IDLE:
      RCLCPP_INFO(this->get_logger(), "State: IDLE");
      break;
      
    case ExplorationState::DETECTING_ENVIRONMENT:
      RCLCPP_INFO(this->get_logger(), 
        "State: DETECTING_ENVIRONMENT");
      break;
      
    case ExplorationState::EXPLORING:
      RCLCPP_INFO(this->get_logger(), "State: EXPLORING");
      break;
      
    case ExplorationState::STUCK:
      RCLCPP_WARN(this->get_logger(), "State: STUCK");
      break;
      
    case ExplorationState::RECOVERY:
      RCLCPP_WARN(this->get_logger(), 
        "State: RECOVERY - Executing recovery behavior");
      executeRecoveryBehavior();
      break;
      
    default:
      break;
  }
}
\end{lstlisting}

\chapter{Système de Sécurité et Contrôle d'Arrêt d'Urgence}

\section{Vue d'ensemble}

La sécurité est un aspect critique du système robotique. Le système d'arrêt d'urgence (Emergency Stop Controller) fonctionne en parallèle avec les autres composants et peut arrêter le robot à tout moment si un danger est détecté.

\section{Architecture du Contrôleur d'Arrêt d'Urgence}

\begin{lstlisting}
EmergencyStopController::EmergencyStopController(
  const rclcpp::NodeOptions & options)
: Node("emergency_stop_controller", options),
  current_safety_state_(SafetyState::SAFE),
  is_in_recovery_(false),
  emergency_stop_active_(false)
{
  // Declare safety parameters
  this->declare_parameter("emergency_stop_distance", 0.15);
  this->declare_parameter("warning_distance", 0.25);
  this->declare_parameter("inflation_radius", 0.60);
  this->declare_parameter("collision_detection_distance", 0.10);
  this->declare_parameter("safety_check_frequency", 50.0);
  this->declare_parameter("max_linear_velocity", 0.4);
  this->declare_parameter("max_angular_velocity", 1.0);
  this->declare_parameter("recovery_backup_distance", 0.3);
  this->declare_parameter("recovery_rotation_angle", 1.57);
  this->declare_parameter("recovery_timeout", 10.0);
  
  // Initialize TF2
  tf_buffer_ = std::make_unique<tf2_ros::Buffer>(
    this->get_clock());
  tf_listener_ = std::make_shared<tf2_ros::TransformListener>(
    *tf_buffer_);
}
\end{lstlisting}

\section{Abonnements aux Capteurs}

Le contrôleur s'abonne aux capteurs critiques:

\begin{lstlisting}
// Subscribe to laser scan data (high priority)
laser_sub_ = this->create_subscription<sensor_msgs::msg::LaserScan>(
  "/scan", rclcpp::SensorDataQoS(),
  std::bind(&EmergencyStopController::laserScanCallback, this, 
    std::placeholders::_1));

// Subscribe to velocity commands
cmd_vel_sub_ = this->create_subscription<geometry_msgs::msg::Twist>(
  "/cmd_vel", 10,
  std::bind(&EmergencyStopController::cmdVelCallback, this, 
    std::placeholders::_1));

// Subscribe to odometry
odom_sub_ = this->create_subscription<nav_msgs::msg::Odometry>(
  "/odom", 10,
  std::bind(&EmergencyStopController::odomCallback, this, 
    std::placeholders::_1));
\end{lstlisting}

\section{Détection des États de Sécurité}

Le système détecte quatre états de sécurité distincts:

\begin{lstlisting}
enum class SafetyState {
  SAFE,                    // No obstacles in danger zone
  WARNING,                 // Obstacles detected at warning distance
  EMERGENCY_STOP,          // Obstacles at emergency stop distance
  COLLISION_DETECTED,      // Collision already occurred
  RECOVERY_NEEDED          // Robot stuck, needs recovery
};
\end{lstlisting}

\section{Analyse de Sécurité}

Le contrôleur analyse continuellement les données des capteurs:

\begin{lstlisting}
void EmergencyStopController::laserScanCallback(
  const sensor_msgs::msg::LaserScan::SharedPtr msg)
{
  std::lock_guard<std::mutex> lock(safety_mutex_);
  latest_scan_ = msg;
  
  // Update safety sectors based on laser data
  updateSafetySectors(msg);
  
  // Analyze current safety state
  current_safety_state_ = analyzeSafetyState(msg);
}

SafetyState EmergencyStopController::analyzeSafetyState(
  const sensor_msgs::msg::LaserScan::SharedPtr msg)
{
  // Find minimum distance to obstacles
  double min_distance = std::numeric_limits<double>::infinity();
  
  for (size_t i = 0; i < msg->ranges.size(); ++i) {
    double range = msg->ranges[i];
    
    if (range > msg->range_min && range < msg->range_max && 
        std::isfinite(range)) {
      min_distance = std::min(min_distance, static_cast<double>(range));
    }
  }
  
  // Determine safety state
  if (min_distance < collision_detection_distance_) {
    return SafetyState::COLLISION_DETECTED;
  } else if (min_distance < emergency_stop_distance_) {
    return SafetyState::EMERGENCY_STOP;
  } else if (min_distance < warning_distance_) {
    return SafetyState::WARNING;
  }
  
  return SafetyState::SAFE;
}
\end{lstlisting}

\section{Gestion des Différents États}

\begin{lstlisting}
void EmergencyStopController::safetyTimerCallback()
{
  std::lock_guard<std::mutex> lock(safety_mutex_);
  
  if (!latest_scan_) {
    return;
  }

  publishSafetyStatus();

  // Handle different safety states
  switch (current_safety_state_) {
    case SafetyState::SAFE:
      emergency_stop_active_ = false;
      last_safe_time_ = this->now();
      // Publish the original velocity command
      publishSafeVelocity(latest_cmd_vel_);
      break;
      
    case SafetyState::WARNING:
      emergency_stop_active_ = false;
      {
        // Reduce velocity in warning zone
        auto reduced_cmd = latest_cmd_vel_;
        reduced_cmd.linear.x *= 0.5;   // 50% speed reduction
        reduced_cmd.angular.z *= 0.7;  // 30% angular reduction
        publishSafeVelocity(reduced_cmd);
      }
      break;
      
    case SafetyState::EMERGENCY_STOP:
    case SafetyState::COLLISION_DETECTED:
      // Execute immediate emergency stop
      executeEmergencyStop();
      break;
      
    case SafetyState::RECOVERY_NEEDED:
      if (!is_in_recovery_) {
        executeRecoveryBehavior();
      }
      break;
  }
}
\end{lstlisting}

\section{Exécution du Comportement de Récupération}

Quand le robot se retrouve bloqué, un comportement de récupération automatique est exécuté:

\begin{lstlisting}
void EmergencyStopController::executeRecoveryBehavior()
{
  is_in_recovery_ = true;
  RCLCPP_WARN(this->get_logger(), 
    "Executing recovery behavior - backing up and rotating");
  
  auto recovery_start = std::chrono::steady_clock::now();
  
  // Phase 1: Backup movement
  geometry_msgs::msg::Twist backup_cmd;
  backup_cmd.linear.x = -0.2;  // Backward movement
  
  // Phase 2: Rotation to escape
  geometry_msgs::msg::Twist rotation_cmd;
  rotation_cmd.angular.z = recovery_rotation_speed_;
  
  // Execute recovery sequence
  while (std::chrono::steady_clock::now() - recovery_start < 
         std::chrono::seconds(static_cast<int>(recovery_timeout_))) {
    publishSafeVelocity(backup_cmd);
    std::this_thread::sleep_for(std::chrono::milliseconds(100));
  }
  
  is_in_recovery_ = false;
  RCLCPP_INFO(this->get_logger(), "Recovery behavior completed");
}
\end{lstlisting}

\chapter{Navigation et Suivi de Mur}

\section{Intégration Navigation2}

Navigation2 fournit une pile de navigation complète pour le calcul de trajectoires sans collision. Notre système l'intègre de manière transparente:

\begin{lstlisting}
// Navigation2 configuration topics
// /plan - computed path
// /local_costmap/costmap - local costmap
// /global_costmap/costmap - global costmap
// /tf - coordinate transforms
\end{lstlisting}

\section{Stratégie de Suivi de Mur}

Le Wall Follower Node implémente une stratégie de navigation alternative pour les environnements complexes:

\begin{lstlisting}
WallFollower::WallFollower(const rclcpp::NodeOptions & options)
: Node("wall_follower", options),
  current_state_(WallFollowingState::SEARCHING_FOR_WALL)
{
  // Declare wall following parameters
  this->declare_parameter("target_wall_distance", 0.4);
  this->declare_parameter("max_wall_detection_range", 2.5);
  this->declare_parameter("min_wall_length", 0.3);
  this->declare_parameter("wall_following_speed", 0.3);
  this->declare_parameter("angular_gain", 2.0);
  this->declare_parameter("distance_gain", 1.5);
  this->declare_parameter("wall_confidence_threshold", 0.7);
  
  // Initialize laser and odometry subscriptions
  laser_sub_ = this->create_subscription<sensor_msgs::msg::LaserScan>(
    "/scan", 10, 
    std::bind(&WallFollower::laserScanCallback, this, 
      std::placeholders::_1));
  
  odom_sub_ = this->create_subscription<nav_msgs::msg::Odometry>(
    "/odom", 10, 
    std::bind(&WallFollower::odomCallback, this, 
      std::placeholders::_1));
  
  // Velocity publisher
  cmd_vel_pub_ = this->create_publisher<geometry_msgs::msg::Twist>(
    "/cmd_vel", 10);
}
\end{lstlisting}

\section{Détection de Mur}

Le système utilise une analyse géométrique des données de scans laser:

\begin{lstlisting}
std::vector<Wall> WallFollower::detectWalls(
  const sensor_msgs::msg::LaserScan::SharedPtr msg)
{
  std::vector<Wall> detected_walls;
  std::vector<cv::Point2f> wall_points;
  
  // Convert laser scan to points
  for (size_t i = 0; i < msg->ranges.size(); ++i) {
    float range = msg->ranges[i];
    
    if (range > msg->range_min && range < msg->range_max && 
        std::isfinite(range)) {
      float angle = msg->angle_min + i * msg->angle_increment;
      
      cv::Point2f point(
        range * std::cos(angle),
        range * std::sin(angle)
      );
      wall_points.push_back(point);
    }
  }
  
  // Fit lines to detect walls
  // Use RANSAC or similar algorithm
  
  return detected_walls;
}
\end{lstlisting}

\chapter{Système de Lancement et Orchestration}

\section{Architecture de Lancement}

Le système utilise des fichiers de lancement Python ROS2 pour orchestrer le démarrage de tous les composants:

\begin{lstlisting}[language=Python]
def generate_launch_description():
    # Package directories
    nav2_bringup_dir = get_package_share_directory('nav2_bringup')
    turtle_bot_nav_dir = get_package_share_directory(
      'turtle_bot_navigation')
    turtle_bot_exp_dir = get_package_share_directory(
      'turtle_bot_exploration')
    
    # Declare launch arguments
    declare_use_sim_time_cmd = DeclareLaunchArgument(
        'use_sim_time',
        default_value='true',
        description='Use simulation (Gazebo) clock')
    
    # Navigation2 bringup
    nav2_bringup_cmd = IncludeLaunchDescription(
        PythonLaunchDescriptionSource(
            os.path.join(nav2_bringup_dir, 'launch', 
              'bringup_launch.py')),
        launch_arguments={
            'map': map_yaml_file,
            'use_sim_time': use_sim_time,
            'params_file': params_file,
        }.items())
    
    # Frontier detection node
    frontier_detection_node = Node(
        package='turtle_bot_exploration',
        executable='frontier_detection_node',
        name='frontier_detection_node',
        parameters=[
            {'use_sim_time': use_sim_time},
            {'frontier_min_size': 15},
            {'frontier_search_radius': 8.0},
        ])
    
    # Exploration coordinator node
    exploration_coordinator_node = Node(
        package='turtle_bot_exploration',
        executable='exploration_coordinator_node',
        parameters=[
            {'use_sim_time': use_sim_time},
            {'exploration_timeout': 1800.0},
            {'auto_start_exploration': True},
        ])
    
    # Create launch description
    ld = LaunchDescription()
    ld.add_action(declare_use_sim_time_cmd)
    ld.add_action(nav2_bringup_cmd)
    ld.add_action(frontier_detection_node)
    ld.add_action(exploration_coordinator_node)
    
    return ld
\end{lstlisting}

\section{Dépendances et Topologie}

La topologie de lancement définit les dépendances entre les nœuds. Chaque nœud attend la disponibilité de certains topics avant de commencer son traitement.

\chapter{Messages ROS2 et Interfaces de Données}

\section{Types de Messages Utilisés}

\subsection{nav\_msgs/OccupancyGrid}

Utilisé pour représenter la carte 2D:

\begin{lstlisting}
std_msgs/Header header
nav_msgs/MapMetaData info
int8[] data  // -1: unknown, 0: free, 1-100: occupied
\end{lstlisting}

\subsection{geometry\_msgs/Twist}

Utilisé pour les commandes de vitesse:

\begin{lstlisting}
geometry_msgs/Vector3 linear
  float64 x     // forward velocity
  float64 y     // lateral velocity
  float64 z     // vertical velocity
geometry_msgs/Vector3 angular
  float64 x     // roll velocity
  float64 y     // pitch velocity
  float64 z     // yaw velocity
\end{lstlisting}

\subsection{visualization\_msgs/Marker et MarkerArray}

Utilisés pour la visualisation en RViz:

\begin{lstlisting}
std_msgs/Header header
string ns
int32 id
int32 type         // SPHERE, CUBE, ARROW, etc.
geometry_msgs/Pose pose
geometry_msgs/Vector3 scale
std_msgs/ColorRGBA color
float32 lifetime
\end{lstlisting}

\subsection{sensor\_msgs/LaserScan}

Données brutes du capteur laser:

\begin{lstlisting}
std_msgs/Header header
float32 angle_min
float32 angle_max
float32 angle_increment
float32 time_increment
float32 scan_time
float32 range_min
float32 range_max
float32[] ranges
float32[] intensities
\end{lstlisting}

\chapter{Paramétrage et Configuration}

\section{Système de Paramètres ROS2}

ROS2 fournit un serveur de paramètres centralisé permettant une configuration dynamique:

\begin{lstlisting}
// Set parameters from command line
ros2 param set /frontier_detector min_frontier_size 20.0

// Get parameters from node
auto min_size = this->get_parameter("min_frontier_size");

// List all parameters
ros2 param list

// Load parameters from YAML file
ros2 run robot_package node_name --ros-args --params-file config.yaml
\end{lstlisting}

\section{Fichiers de Configuration YAML}

Les paramètres sont généralement stockés dans des fichiers YAML:

\begin{lstlisting}
# navigation_params.yaml
frontier_detector:
  ros__parameters:
    min_frontier_size: 15.0
    max_frontier_distance: 10.0
    frontier_travel_point_distance: 0.8

exploration_coordinator:
  ros__parameters:
    exploration_timeout: 1800.0
    stuck_timeout: 8.0
    goal_tolerance: 0.3
    enable_dynamic_tolerance: true
    environmental_adaptation: true

emergency_stop_controller:
  ros__parameters:
    emergency_stop_distance: 0.15
    warning_distance: 0.25
    safety_check_frequency: 50.0
\end{lstlisting}

\chapter{Défis Techniques et Solutions}

\section{Synchronisation Distribuée}

\subsection{Défi}

Dans un système distribué, les nœuds peuvent recevoir des messages avec des horodatages différents, ce qui crée une incohérence temporelle.

\subsection{Solution}

Utilisation des transformations TF2 pour aligner temporellement les données:

\begin{lstlisting}
// Wait for transform and apply it
geometry_msgs::msg::PointStamped point_in;
point_in.header.frame_id = "laser_frame";
point_in.point.x = 1.0;

geometry_msgs::msg::PointStamped point_out;
try {
  tf_buffer_->transform(point_in, point_out, "base_link");
} catch (tf2::TransformException & ex) {
  RCLCPP_ERROR(this->get_logger(), "%s", ex.what());
}
\end{lstlisting}

\section{Gestion de la Concurrence}

\subsection{Défi}

Plusieurs callbacks peuvent accéder aux mêmes données simultanément.

\subsection{Solution}

Utilisation de mécanismes de synchronisation (mutex):

\begin{lstlisting}
private:
  std::mutex safety_mutex_;
  
void laserScanCallback(...) {
  std::lock_guard<std::mutex> lock(safety_mutex_);
  latest_scan_ = msg;
}

void safetyTimerCallback() {
  std::lock_guard<std::mutex> lock(safety_mutex_);
  // Safe access to latest_scan_
}
\end{lstlisting}

\section{Latence et Temps Réel}

\subsection{Défi}

Les délais de communication et de traitement peuvent compromettre la sécurité.

\subsection{Solution}

Implémentation de timers dédiés à haute fréquence pour les tâches critiques:

\begin{lstlisting}
// Safety timer at 50Hz
auto safety_timer_period = std::chrono::milliseconds(
  static_cast<int>(1000.0 / 50.0));
safety_timer_ = this->create_wall_timer(
  safety_timer_period, 
  std::bind(&EmergencyStopController::safetyTimerCallback, this));
\end{lstlisting}

\section{Transformation de Coordonnées}

\subsection{Défi}

Le robot doit travailler dans plusieurs systèmes de coordonnées (base\_link, map, laser, etc.).

\subsection{Solution}

Utilisation de TF2 pour maintenir une arborescence de transformations:

\begin{lstlisting}
// Base-link to map transform
tf2::Stamped<tf2::Transform> base_to_map;
try {
  tf_buffer_->lookupTransform("map", "base_link", 
    tf2::TimePointZero, base_to_map);
  
  auto robot_x = base_to_map.getOrigin().x();
  auto robot_y = base_to_map.getOrigin().y();
} catch (tf2::TransformException & ex) {
  RCLCPP_ERROR(this->get_logger(), "%s", ex.what());
}
\end{lstlisting}

\chapter{Performance et Métriques}

\section{Métriques Clés}

\begin{itemize}
  \item \textbf{Couverture de carte:} Pourcentage de l'environnement cartographié
  \item \textbf{Efficacité d'exploration:} Ratio entre surface explorée et distance parcourue
  \item \textbf{Latence de détection:} Temps entre la détection d'un obstacle et l'arrêt
  \item \textbf{Fréquence de contrôle:} Fréquence de traitement des cycles de contrôle
  \item \textbf{Fiabilité SLAM:} Probabilité de succès de localisation
\end{itemize}

\section{Optimisations Implémentées}

\subsection{Sélection Intelligente de Frontière}

Le système utilise une fonction de score multi-critères pour sélectionner les frontières optimales, équilibrant distance, gain informatif et faisabilité.

\subsection{Filtrage Dynamique de Tolérance}

Les tolérances de navigation s'ajustent dynamiquement basées sur la vitesse et les conditions de l'environnement.

\subsection{Adaptation Environnementale}

Le système apprend des patterns de l'environnement pour optimiser les futures explorations.

\chapter{Intégration avec l'Application Mobile}

\section{Architecture Distribuée Client-Serveur}

Le système ROS2 peut être accédé par une application mobile via une couche de communication réseau. Les topics ROS2 sont exposés via une API REST ou WebSocket.

\section{Topics Exposés}

\begin{itemize}
  \item \textbf{/exploration/status:} État actuel de l'exploration
  \item \textbf{/map:} Carte en cours de construction
  \item \textbf{/tf:} Transformations de coordonnées
  \item \textbf{/robot\_pose:} Position actuelle du robot
  \item \textbf{/cmd\_vel:} Commandes de vitesse (pour contrôle manuel)
\end{itemize}

\section{Services ROS2 pour Contrôle Mobile}

\begin{lstlisting}
// Service: StartExploration
// Demande au système de commencer l'exploration

// Service: StopExploration
// Arrête l'exploration en cours

// Service: EmergencyStop
// Déclenche l'arrêt d'urgence

// Service: GetRobotStatus
// Retourne l'état actuel du robot
\end{lstlisting}

\chapter{Conclusion et Perspectives Futures}

\section{Accomplissements}

Ce projet a avec succès implémenté un système complet de navigation autonome et d'exploration basé sur ROS2, incluant:

\begin{enumerate}
  \item Un système robuste de détection des frontières avec optimisation multi-critères
  \item Un coordinateur d'exploration sophistiqué avec machine à états complète
  \item Un système de sécurité multicouche avec arrêt d'urgence et récupération
  \item Une intégration seamless avec Navigation2 et SLAM
  \item Un framework de visualisation en temps réel
  \item Une architecture extensible pour l'intégration avec des interfaces utilisateur
\end{enumerate}

\section{Améliorations Futures}

\begin{enumerate}
  \item \textbf{Apprentissage par renforcement:} Utiliser des techniques d'IA pour optimiser la sélection de frontière
  \item \textbf{Navigation multi-robot:} Étendre le système pour coordonner plusieurs robots
  \item \textbf{Exploration avec contraintes:} Intégrer des contraintes énergétiques et temporelles
  \item \textbf{Cartographie sémantique:} Ajouter la reconnaissance d'objets à la cartographie
  \item \textbf{Adaptation temps réel:} Améliorer les mécanismes d'adaptation aux changements d'environnement
\end{enumerate}

\section{Applicabilité Industrielle}

Le système développé a des applications pratiques dans:

\begin{itemize}
  \item Recherche et sauvetage (SAR)
  \item Inspection d'environnements dangereux
  \item Cartographie et surveillance
  \item Nettoyage autonome
  \item Livraison intérieure
\end{itemize}

\newpage

\begin{thebibliography}{99}

\bibitem{ros2_design} Maruyama, Y., Kato, S., \& Azumi, T. (2016). Exploring the capabilities of ROS2. In 13th International Conference on Embedded Software, EMSOFT 2016.

\bibitem{frontier_exploration} Yamauchi, B. (1997). A frontier-based approach for autonomous exploration. In Proceedings 1997 IEEE International Symposium on Computational Intelligence in Robotics and Automation.

\bibitem{nav2} Macenski, S., Foote, T., Gerkey, B., Lalancette, C., \& Woodall, W. (2021). Robot Operating System 2: Design, architecture, and uses in the wild. Science and Technology, 2.

\bibitem{costmap_layers} Marder-Eppstein, E., Berger, E., Foote, T., Gerkey, B., \& Konolige, K. (2010). The office marathon: Robust navigation in an indoor office environment. In Robotics and Automation, 2010.

\bibitem{dwa_local_planner} Fox, D., Burgard, W., \& Thrun, S. (1997). The dynamic window approach to collision avoidance. In Robotics and Automation, 1997.

\bibitem{slam_toolbox} Macenski, S., Goins, G., \& Gerkey, B. (2021). Comparison of Cost Map Inflation Models in Navigation2. In 2021 IEEE International Conference on Robotics and Automation.

\bibitem{tf2_transforms} Foote, T., \& Geyer, C. (2013). TF2. In IEEE Robotics \& Automation Magazine.

\bibitem{qos_policy} Data Distribution Service. (2015). Interoperability Wire Protocol Specification. Object Management Group.

\bibitem{mobile_robotics} Siegwart, R., Nourbakhsh, I. R., \& Scaramuzza, D. (2011). Introduction to autonomous mobile robots. MIT press.

\bibitem{path_planning} LaValle, S. M. (2006). Planning algorithms. Cambridge university press.

\end{thebibliography}

\appendix

\chapter{Code Complet des Nœuds Clés}

\section{Frontier Detector - Structure Complète}

Les sources complètes sont disponibles dans le package turtle\_bot\_exploration. Le nœud frontier detector comprend:

\begin{itemize}
  \item Detection itérative par BFS (Breadth-First Search)
  \item Clustering de cellules frontière
  \item Calcul de gain informatif
  \item Filtrage par distance et taille
  \item Vérification d'accessibilité par pathfinding
\end{itemize}

\section{Exploration Coordinator - Structure Complète}

Le coordinateur d'exploration contient:

\begin{itemize}
  \item Machine à états avec 8 états distincts
  \item Sélection intelligente de frontière multi-critères
  \item Gestion d'actions Navigation2
  \item Détection d'état bloqué et récupération
  \item Monitoring de performance et métriques
  \item Adaptation dynamique de paramètres
\end{itemize}

\section{Emergency Stop Controller - Structure Complète}

Le contrôleur de sécurité fournit:

\begin{itemize}
  \item Analyse de sécurité à 50Hz
  \item Gestion des états de sécurité (SAFE, WARNING, EMERGENCY)
  \item Modulation de vitesse adaptative
  \item Comportement de récupération automatique
  \item Logging détaillé et status reporting
\end{itemize}

\chapter{Configuration Complète des Paramètres}

\begin{lstlisting}
# Configuration complète pour tous les nœuds

# Frontier Detector Configuration
frontier_detector:
  ros__parameters:
    min_frontier_size: 15.0
    max_frontier_distance: 10.0
    frontier_travel_point_distance: 0.8
    frontier_blacklist_radius: 1.2
    frontier_search_radius: 8.0
    frontier_min_distance: 1.0
    information_radius: 1.0
    robot_base_frame: "base_link"
    global_frame: "map"

# Exploration Coordinator Configuration  
exploration_coordinator:
  ros__parameters:
    exploration_timeout: 1800.0
    stuck_timeout: 8.0
    min_map_coverage_to_start: 0.1
    goal_tolerance: 0.3
    stuck_distance_threshold: 0.15
    goal_reached_tolerance: 1.2
    goal_timeout: 45.0
    force_new_goal_distance: 1.5
    recovery_rotation_speed: 0.5
    recovery_rotation_duration: 12.0
    enable_dynamic_tolerance: true
    velocity_based_tolerance: true
    environmental_adaptation: true
    performance_monitoring: true
    adaptive_parameter_tuning: true
    auto_start_exploration: true

# Emergency Stop Controller Configuration
emergency_stop_controller:
  ros__parameters:
    emergency_stop_distance: 0.15
    warning_distance: 0.25
    inflation_radius: 0.60
    collision_detection_distance: 0.10
    safety_check_frequency: 50.0
    max_linear_velocity: 0.4
    max_angular_velocity: 1.0
    recovery_backup_distance: 0.3
    recovery_rotation_angle: 1.57
    recovery_timeout: 10.0
    base_frame: "base_link"
    laser_frame: "laser"

# Navigation2 Parameters
nav2_params:
  planner_server:
    ros__parameters:
      planner_plugins: ["GridBased"]
      GridBased:
        plugin: "nav2_navfn_planner/NavfnPlanner"
  
  controller_server:
    ros__parameters:
      controller_plugins: ["FollowPath"]
      FollowPath:
        plugin: "nav2_regulated_pure_pursuit_controller/RegulatedPurePursuitController"
        desired_linear_vel: 0.2
        max_angular_vel: 1.0

  local_costmap:
    ros__parameters:
      update_frequency: 5.0
      publish_frequency: 2.0
      global_frame: "odom"
      robot_base_frame: "base_link"
      
  global_costmap:
    ros__parameters:
      update_frequency: 1.0
      publish_frequency: 1.0
      global_frame: "map"
      robot_base_frame: "base_link"
\end{lstlisting}

\chapter{Outils de Diagnostic et Monitoring}

\section{Commandes ROS2 Utiles}

\begin{lstlisting}[language=bash]
# Liste tous les nœuds en cours d'exécution
ros2 node list

# Liste tous les topics
ros2 topic list

# Affiche les détails d'un topic
ros2 topic info /exploration/status

# Écoute les messages d'un topic
ros2 topic echo /exploration/status

# Affiche les statistiques de topic
ros2 topic hz /exploration/status

# Liste tous les services disponibles
ros2 service list

# Appelle un service
ros2 service call /start_exploration std_srvs/srv/Empty

# Affiche les paramètres d'un nœud
ros2 param list /exploration_coordinator

# Définit un paramètre
ros2 param set /frontier_detector min_frontier_size 20.0

# Lance RViz pour la visualisation
ros2 run rviz2 rviz2 -d exploration.rviz

# Affiche le graph de tf
ros2 run tf2_tools view_frames
\end{lstlisting}

\section{Debugging Avancé}

Pour le debugging détaillé, activer le logging de debug:

\begin{lstlisting}[language=bash]
# Logging au niveau DEBUG pour un nœud
ros2 run turtle_bot_exploration frontier_detection_node \
  --ros-args --log-level frontier_detector:=DEBUG

# Enregistrement bag pour analyse offline
ros2 bag record /map /exploration/status /tf

# Lecture du bag enregistré
ros2 bag play rosbag2_2024_01_15-14_23_45/rosbag2_2024_01_15-14_23_45_0.db3
\end{lstlisting}

\end{document}
